% !PRINTED 86
\weberpage{86}

%<<CONT
worin die Funktionen $f(\omega)$, $f_1(\omega)$, $f_2(\omega)$
folgendermaßen definiert sind:
\weq{11}{
\begin{aligned}
f(\omega) &= q^{-\frac{1}{24}}\prod_{1,\infty}^{\nu}(1+q^{2\nu-1}),\\
f_1(\omega) &= q^{-\frac{1}{24}}\prod_{1,\infty}^{\nu}(1-q^{2\nu-1}),\\
f_2(\omega) &= \sqrt{2}\,q^{\frac{1}{12}}\prod_{1,\infty}^{\nu}(1+q^{2\nu}).
\end{aligned}
}

Die Funktionen $\eta(\omega)$, $f(\omega)$, $f_1(\omega)$,
$f_2(\omega)$ werden in unseren späteren Betrachtungen eine
wichtige Rolle spielen. Aus § 21, (14) ergibt sich nach (10)
die Relation
\weq{12}{f(\omega)^8 = f_1(\omega)^8 + f_2(\omega)^8,}
und aus § 23 (5) nach (9)
\weq{13}{f(\omega)f_1(\omega)f_2(\omega) = \sqrt{2}.}
Die letzte Formel ergibt sich auch aus (11) mit Benutzung der
identischen Relation:
\[ \prod (1+q^\nu)(1-q^{2\nu-1}) = \prod \frac{(1+q^\nu)(1-q^\nu)}{1-q^{2\nu}} = 1. \]

\begin{center}
\textbf{§ 25. Darstellung der $\vartheta$-Funktionen durch
unendliche Reihen.}
\end{center}

Der zweite Weg, um zur Darstellung der $\vartheta$-Funktionen zu
gelangen, besteht darin, daß man eine den Fundamentalgleichungen
genügende konvergente unendliche Reihe zu bilden sucht.

Bemerken wir zunächst, daß wegen der Bedingung
\[ \vartheta_{00}(v+1) = \vartheta_{00}(v) \]
die Funktion $\vartheta_{00}(v)$ als eindeutige Funktion von
$e^{2\pi i v}$ angesehen werden kann, und setzen demgemäß
\[ \vartheta_{00}(v) = \sum_{-\infty,\infty}^{\nu} A_\nu e^{2\pi i v \nu}, \]
so ergibt die Differentialgleichung § 20, (4):
\[ \frac{\partial^2 \vartheta}{\partial v^2} - 4\pi i \frac{\partial \vartheta}{\partial \omega} = 0 \]
für $A_\nu$ die Bedingung
\[ \frac{dA_\nu}{d\omega} = \pi i \nu^2 A_\nu, \qquad A_\nu = c_\nu e^{\pi i \omega \nu^2} = c_\nu q^{\nu^2}, \]
worin $c_\nu$ in bezug auf $\omega$ konstant ist. Es wird also
\[ \vartheta_{00}(v) = \sum c_\nu q^{\nu^2} c^{2\pi i v \nu}, \]
% [sic] Printed c as the exponential base; mathematically expected e.
%CONT>>
